\documentclass[10pt,a4paper]{amsart}
\usepackage[margin=19mm]{geometry}
\usepackage{amsmath,amssymb,mathtools}
\usepackage[scr=rsfs,cal=euler]{mathalfa}
\usepackage{xfrac}
\usepackage[unicode,hidelinks,bookmarks=false]{hyperref}
\newcommand{\ie}{i.e.\ }
\newcommand{\cf}{cf.\ }
\newcommand{\resp}{respectively\ }
\newcommand{\Subsection}{\subsection}
\providecommand{\nobreakdash}{\nobreak\hbox{-}}
\setlength{\emergencystretch}{2em}
\multlinegap0pt
\usepackage{bm}
\usepackage{bbm}
\usepackage{dsfont}
\usepackage{xspace}
\usepackage{cancel}
\usepackage{mathtools}
\usepackage{amsthm}
\makeatletter
\@ifundefined{Theorem}{\newtheorem{Theorem}{Theorem}}{}
\@ifundefined{Proposition}{\newtheorem{Proposition}[Theorem]{Proposition}}{}
\makeatother
\title[Reformulation of $q$-Middle Convolution and Applications]{Reformulation of $q$-Middle Convolution and Applications}
\author{Yumi Arai, Kouichi Takemura}
\date{Source: arXiv:2503.11214v3 (CC BY 4.0)}
\begin{document}
\maketitle

\noindent\emph{Source and licence.} Reformulation of $q$-Middle Convolution and Applications. Version arXiv:2503.11214v3 (CC BY 4.0), distributed under the Creative Commons Attribution 4.0 International licence, as recorded in the arXiv licence metadata for this exact version and shown on the abstract page https://arxiv.org/abs/2503.11214v3. Full rights notice: licenses/arxiv-2503.11214-CC-BY-4.0.txt.\par\smallskip

\noindent\textit{Editorial scope.} Counted unit: the Proposition whose source label is \texttt{prop:qcintconvYx} in the article (source0.tex lines 549--560, source label \texttt{prop:qcintconvYx}), delivered together with the article's own complete original proof of it (source0.tex lines 561--668, one \texttt{proof} environment of 108 lines). No mathematics is written, repaired or re-derived: statement and proof are the article's own text line for line. Required context retained uncounted: none. Every definition, display and label of those lines is retained unchanged. Omitted from this package and not redistributed: 1--548, 669--2576. Nothing inside a retained excerpt is omitted or altered: every retained body is delivered exactly as the article's lines are written, so no omission receipt of any kind accompanies this package. Citations: the retained excerpts carry no citation command at all, so no bibliography entry of the article is transcribed and no reference is redistributed. Reference adaptation: none was needed and none was made; every reference inside the delivered unit resolves inside it, so no unresolved reference or citation is produced. Printed slips kept literal: no printed word, symbol, hypothesis or number of the counted statement or of its proof was altered. Layout and class: the article's own class is \texttt{sigma} with packages including ifpdf, tikz; the wrapper is the lane's pinned \texttt{amsart} editorial wrapper at 10pt on A4, which restates the theorem-like environments the retained text needs and adds the article's own macro definitions that the retained text needs (none), with extra wrapper packages (bm, bbm, dsfont, xspace, cancel, mathtools). The delivered numbering is the unit's own. Assets: no retained span contains a figure environment, a graphics-inclusion command, a PDF-page inclusion, a table or a TikZ picture; no image file is copied into this package, the package declares no asset dependency and no image file is redistributed with it. Licence: version arXiv:2503.11214v3 of this article is distributed under the Creative Commons Attribution 4.0 International licence, as recorded in the arXiv licence metadata for this exact version and shown on the abstract page https://arxiv.org/abs/2503.11214v3. A full rights notice is delivered at licenses/arxiv-2503.11214-CC-BY-4.0.txt. Characters: every retained excerpt is pure ASCII, so the delivered engine, XeLaTeX with its default Latin Modern text font, reports no missing character.
\begin{Proposition} \label{prop:qcintconvYx}
Set $b_0=0$.
Let $Y(x)$ be a solution to
\begin{equation}
\frac{Y(q x) - Y(x )}{-x} = \Biggl[ \sum^{N}_{i = 0}\frac{B_{i}}{x -b_{i}} \Biggr] Y(x).
\label{eq:YqxBYx00}
\end{equation}
\begin{itemize}\itemsep=0pt
\item[$(i)$] If the absolute value of any eigenvalue of $I_m - B_0$ is strictly less than $1$, then there exist~${\varepsilon _1 , C_1 \in \mathbb{R}_{>0}} $ and $M_1 \in \mathbb{Z} $ such that $|[Y(s )]_k | \leq C_1 |s |^{\varepsilon _1} $ for any $s \in \{ q^{n}\xi \mid n \geq M_1 ,\, n \in \mathbb{Z} \} $ and $k=1,\dots ,m$.
\item[$(ii)$] If the absolute value of any eigenvalue of $I_m - B_0 -B_1 - \dots - B_N$ is strictly more than~$\big|q^{\lambda }\big|$, then there exist $\varepsilon _2 , C_2 \in \mathbb{R}_{>0} $ and $M_2 \in \mathbb{Z} $ such that $|[Y(s )]_k | \leq C_2 \big|s ^{\lambda }\big| |s | ^{ - \varepsilon _2}$ for any~${s \in \{ q^{n}\xi \mid n \leq -M_2 ,\, n \in \mathbb{Z} \} }$ and $k=1,\dots ,m$.
\end{itemize}
\end{Proposition}

\begin{proof}
Write equation \eqref{eq:YqxBYx00} as
$
Y(q x) = B(x) Y(x) $.
Then, we have $B(x) \to I_m -B_0$ as $x \to 0$ and $B(x) \to I_m -B_0 -B_1 - \dots - B_N$ as $x \to \infty$.

We show (i).
Let $e_1, e_2, \dots ,e_m$ be the eigenvalues of the matrix $I_m- B_0 $.
By the assumption, we have $\max \{ |e_1|, |e_2|, \dots , |e_m| \} <1 $.
There exists $\varepsilon _1 >0$ such that $|e_k | <|q|^{\varepsilon _1} $ for any $k \in \{1, \dots ,m \}$.
Set $ \varepsilon = |q|^{\varepsilon _1} - \max \{ |e_1|, |e_2|, \dots , |e_m| \} (>0)$.
There exist numbers $p_j \in \{ 0, \varepsilon /2 \}$ $(j=1, \dots , m-1)$ and an invertible matrix $Q_0$ such that
\begin{equation*}
 Q_0^{-1} ( I_m -B_0 ) Q_0 =
\begin{pmatrix}
e_1 & p_1 & \cdots & 0 & 0 \\
0 & e_2 & \ddots & 0 & 0 \\
\vdots & \ddots & \ddots & \ddots & \vdots \\
0 & 0 & \ddots & e_{m-1} & p_{m-1} \\
0 & 0 & \cdots & 0 & e_{m}
\end{pmatrix}.
\end{equation*}
Set $p_m=0$ for later use.
The numbers $p_1, \dots ,p_{m-1}$ are determined by the Jordan normal form of $ I_m -B_0 $.
If $I_m- B_0 $ is diagonalizable, then $p_1 = \dots = p_{m-1} =0$.
Write
\begin{equation} \label{eq:Q0-BxQ0}
 Q_0^{-1} B(x) Q_0 = Q_0^{-1} ( I_m -B_0 ) Q_0 + \begin{pmatrix}
b_{1,1} (x) & \cdots & b_{1,m} (x) \\
\vdots & & \vdots \\
 b_{m,1} (x) & \cdots & b_{m,m} (x)
\end{pmatrix}.
\end{equation}
Since $Q_0^{-1} B(x) Q_0 \to Q_0^{-1} ( I_m -B_0 ) Q_0 $ as $x \to 0$, there exists $\delta >0$ such that $|b_{i,j}(x)| < \varepsilon /(2m)$ for any $i,j \in \{1, \dots ,m \}$ and any $x$ $(|x| < \delta )$.
Let $Y (x)$ be a solution to $Y (qx) = B(x)Y (x)$.
Then, the function $Z(x) = Q_0^{-1} Y (x)$ satisfies $Z(qx) = Q_0^{-1} B(x) Q_0 Z(x)$.
Let $[Z(s )]_k $ be the $k$-th component of $Z(s ) (\in \mathbb{C} ^{m} )$.
We have
\begin{equation*}
[Z(qx )]_k = e_k [Z(x )]_k + p_k [Z(x )]_{k+1} + b_{k,1}(x) [Z(x )]_1 + \dots + b_{k,m}(x) [Z(x )]_m.
\end{equation*}
We assume that $|x| < \delta $ and $|[Z(x )]_k| < \widetilde{A}$ for any $k \in \{1, \dots ,m \}$.
Then,
\begin{gather}
 | [Z(qx )]_k | \leq |e_k| | [Z(x )]_k| + |p_k | | [Z(x )]_{k+1}| + |b_{k,1}(x)| | [Z(x )]_1| + \dots + |b_{k,m}(x) | |[Z(x )]_m | \nonumber \\
\qquad \leq \widetilde{A} |e_k| + \widetilde{A} \varepsilon /2 + \widetilde{A} m \varepsilon /(2m) \leq \widetilde{A} ( \max \{ |e_1|, |e_2|, \dots , |e_m| \} + \varepsilon ) = |q|^{\varepsilon _1} \widetilde{A} \label{eq:Zqxineq}
\end{gather}
for any $k \in \{1, \dots ,m \}$.
We also have
\begin{equation}
 | [Z(q^n x )]_k | \leq |q|^{n \varepsilon _1} \widetilde{A} \label{eq:ineq}
\end{equation}
for $n \in \mathbb{Z} _{\geq 0} $ and $k \in \{1, \dots ,m \}$.
Let $\xi \in \mathbb{C} \setminus \{0\} $.
There exists an integer $M_0$ such that $| q^n \xi | < \delta $ for any integer $n \geq M_0$.
We take a positive number $C' $ such that $\big| \big[Z\bigl( q^{M_0} \xi \bigr)\big]_k\big| \leq C' $ for any~${k \in \{1, \dots ,m \}}$.
It follows from the inequality \eqref{eq:ineq} that
\[
\big| \big[Z\bigl( q^{M_0 +n} \xi \bigr)\big]_k\big| \leq |q|^{n \varepsilon _1 } C' = \big| q^{M_0 +n} \xi \big|^{\varepsilon _1} \big| q^{M_0 } \xi \big|^{-\varepsilon _1} C'
\]
for any $k \in \{1, \dots ,m \}$ and $n \in \mathbb{Z} _{\geq 0}$.
Since $Y (x) = Q_0 Z(x)$, there exists a positive number $C$ such that $\big| \big[Y\bigl( q^{M_0 +n} \xi \bigr)\big]_k\big| \leq \big|q^{M_0 +n} \xi\big|^{\varepsilon _1} C $ for any $k \in \{1, \dots ,m \}$ and $n \in \mathbb{Z} _{\geq 0}$.
Therefore, we obtain (i).

We show (ii).
Let $e'_1, e'_2, \dots ,e'_m$ be the eigenvalues of the matrix $I_m- B_0 -B_1 - \dots - B_N$.
By the assumption, we have $\min \{ |e'_1|, |e'_2|, \dots , |e'_m| \} > \big|q^{\lambda }\big|$.
Hence, the matrix $I_m -B_0 -B_1 - \dots - B_N$ is invertible and there exists $\varepsilon _2 >0$ such that \smash{$|e'_k |^{-1} < |q | ^{\varepsilon _2} \big|q ^{\lambda }\big|^{-1} $} for any $k \in \{1, \dots ,m \}$.
Set $e_k =1/e'_k $ $(k=1,\dots ,m )$ and \smash{$\varepsilon = |q | ^{\varepsilon _2} \big|q ^{\lambda }\big|^{-1} - \max \{ |e_1|, |e_2|, \dots , |e_m| \} (>0)$}.
There exist numbers $p_j \in \{ 0, \varepsilon /2 \}$ $(j=1, \dots , m-1)$ and an invertible matrix $Q_{\infty }$ such that $Q_{\infty }^{-1} ( I_m -B_0 -B_1 - \dots - B_N ) ^{-1} Q_{\infty }$ is expressed as the right-hand side of equation \eqref{eq:Q0-BxQ0}.
Set $p_m=0$ for later use.
Since $B(x) \to I_m -B_0 -B_1 - \dots - B_N $ as $x \to \infty $, the matrix $B(x)$ is invertible if $|x|$ is sufficiently large.
Write
\begin{align*}
 Q_{\infty }^{-1} B\bigl(q^{-1} x\bigr)^{-1} Q_{\infty } = {}& Q_{\infty }^{-1} ( I_m -B_0 -B_1 - \dots - B_N ) ^{-1} Q_{\infty } \\
& + \begin{pmatrix}
b_{1,1} (x) & \cdots & b_{1,m} (x) \\
\vdots & & \vdots \\
 b_{m,1} (x) & \cdots & b_{m,m} (x)
\end{pmatrix}.
\end{align*}
Since $Q_{\infty }^{-1} B\bigl(q^{-1} x\bigr)^{-1} Q_{\infty } \to Q_{\infty }^{-1} ( I_m -B_0 -B_1 - \dots - B_N ) ^{-1} Q_{\infty } $ as $x \to \infty $, there exists $D >0$ such that $|b_{i,j}(x)| < \varepsilon /(2m)$ for any $i,j \in \{1, \dots ,m \}$ and any $x$ $(|x| > D )$.
Let $Y (x)$ be a solution to $Y (qx) = B(x)Y (x)$.
Then, the function $Z(x) = Q_{\infty }^{-1} Y (x)$ satisfies $Z\bigl(q^{-1}x\bigr) = \smash{Q_{\infty }^{-1} B\bigl(q^{-1}x\bigr) ^{-1} Q_{\infty } Z(x)}$.
We assume that $|x| > D $ and \smash{$|[Z(x )]_k| < \widetilde{A}$} for any $k \in \{1, \dots ,m \}$.
It is shown as equation \eqref{eq:Zqxineq} that \smash{$\big| \big[Z\bigl(q^{-1}x \bigr)\big]_k \big| \leq |q | ^{\varepsilon _2} \big|q ^{\lambda }\big|^{-1} \widetilde{A} $} and
\begin{equation}
 | [Z(q^{-n} x )]_k | \leq |q | ^{n \varepsilon _2} \big|q ^{\lambda }\big|^{-n} \widetilde{A} \label{eq:ineq2}
\end{equation}
for $n \in \mathbb{Z} _{\geq 0} $ and $k \in \{1, \dots ,m \}$.

Let $\xi \in \mathbb{C} \setminus \{0\} $.
There exists an integer $M_{\infty }$ such that $| q^n \xi | < D $ for any integer $n \leq -M_{\infty }$.
We take a positive number $C' $ such that $\big| \big[Z\bigl( q^{-M_{\infty }} \xi \bigr)\big]_k\big| \leq C' $ for any $k \in \{1, \dots ,m \}$.
It follows from the inequality \eqref{eq:ineq2} that
\begin{align*}
\big| \big[Z\bigl( q^{-M_{\infty } -n} \xi \bigr)\big]_k\big| &\leq |q | ^{n \varepsilon _2} \big|q ^{\lambda }\big|^{-n} C'\\
& = \big|\bigl(q^{-M_{\infty } -n } \xi \bigr)^{\lambda } \big| \big|q^{-M_{\infty } -n} \xi \big|^{ - \varepsilon _2 }
  \big|\bigl(q^{-M_{\infty } } \xi \bigr)^{- \lambda } \big| \big|q^{-M_{\infty } } \xi \big|^{ \varepsilon _2 } C' \end{align*}
 for any $k \in \{1, \dots ,m \}$ and $n \in \mathbb{Z} _{\geq 0}$.
Since $Y (x) = Q_{\infty } Z(x)$, we obtain that there exists a~positive number $C$ such that
\[
\big| \big[Y\bigl( q^{-M_{\infty } -n} \xi \bigr)\big]_k\big| \leq \big|\bigl(q^{-M_{\infty } -n } \xi \bigr)^{\lambda } \big|
\cdot \big|q^{-M_{\infty } -n} \xi \big|^{ - \varepsilon _2 } C
\]
 for any $k \in \{1, \dots ,m \}$ and $n \in \mathbb{Z} _{\geq 0}$.
Therefore, we obtain (ii).
\end{proof}

\end{document}
